\section{Error budget, validation, and release criteria}

\subsection{Error decomposition}

The total observable error is decomposed as
\[
 \epsilon_{\mathrm{tot}}
 \lesssim
 \epsilon_{\mathrm{geom}}
 +\epsilon_{\mathrm{basis}}
 +\epsilon_{\mathrm{quad}}
 +\epsilon_{\mathrm{solve}}
 +\epsilon_{\mathrm{macro}}
 +\epsilon_{\mathrm{NN}}
 +\epsilon_{\mathrm{time}}.
\]
The terms denote continuous-geometry approximation, modal truncation,
quadrature, integral solve, rational reduction, neural approximation, and time
integration errors. A release report must not hide one term inside a generic
training loss.

\subsection{Electromagnetic validation}

For held-out scenes and frequency points, report:
\begin{enumerate}
\item port impedance error in physically scaled Frobenius and current-space
      norms;
\item passivity/reciprocity violations before and after the structural decoder;
\item absorbed and outward power closure;
\item true integral-equation residual for CERTIFIED predictions;
\item conductor current and selected near-field probe errors against REFERENCE;
\item loss-channel and continuous spatial-loss errors.
\end{enumerate}

\subsection{Thermal validation}

Report:
\begin{enumerate}
\item source-to-temperature transfer error over the declared Laplace/frequency
      window;
\item full trajectory error at fixed hidden scenes and inputs;
\item peak component and pointwise temperature errors;
\item steady-state residual;
\item energy-balance residual;
\item stability margin of the reduced thermal dynamics.
\end{enumerate}

\subsection{Cross-method validation}

A mesh-free-first teacher must be checked against at least one independent
method on a small benchmark suite. The legacy H(curl) volume solver may serve as
one such oracle after its own convergence is demonstrated; external commercial
or independent open-source solvers are also admissible. Agreement between two
discretizations is evidence, not a proof, and must be accompanied by internal
residual/convergence studies.

\subsection{Release partitions}

The data lifecycle has four disjoint roles:
\begin{description}
\item[training] optimized by gradient descent and active learning;
\item[validation] fixed model-selection scenes;
\item[test] fixed development-wide regression suite;
\item[release] locked scenes, material regimes, poses, frequencies, and
  trajectories not used in architecture or hyperparameter decisions.
\end{description}

\subsection{Fail-closed semantics}

An artifact is production certified only if every required invariant and error
budget passes. Otherwise the artifact may be saved as an experimental model,
but the API must expose its non-certified status. Projection or clipping
corrections are reported explicitly and cannot be used to conceal a poor raw
approximation.
